\subsection{连续线性表示}

设 \(\mathbf{G}\) 为拓扑群，\(\mathbf{E}\) 为局部凸空间，\(U\) 为 \(\mathbf{G}\) 在 \(\mathbf{E}\) 中的线性表示。

定义 1. --- (i) 若对每个 \(s \in G\)，\(U(s)\) 是 E 的连续自同态，并且对每个 \(x \in E\)，从 G 到 E 的映射 \(s \mapsto U(s)x\) 连续，则称 U 分别连续。

\begin{enumerate}
\def\labelenumi{(\roman{enumi})}
\setcounter{enumi}{1}
\item
  若 \(U\) 满足 \((s,x)\mapsto U(s)x\) 是从 \(\mathbf{G}\times \mathbf{E}\) 到 \(\mathbf{E}\) 的连续映射，则称 U 连续。
\item
  若 \(U\) 连续，且由 \(U(s)\)（当 \(s\) 遍历 \(G\) 时）构成的自同态集合等度连续，则称 U 等度连续。
\end{enumerate}

注。--- 1) 说 U 分别连续，是指 \(s \mapsto U(s)\) 是从 G 到连续自同态空间 \(\mathcal{L}(\mathrm{E};\mathrm{E})\) 的连续映射，并在该空间上赋予逐点收敛拓扑。

\begin{enumerate}
\def\labelenumi{\arabic{enumi})}
\setcounter{enumi}{1}
\tightlist
\item
  说 \(U\) 连续，等价于以下三个条件：
\end{enumerate}

\begin{enumerate}
\def\labelenumi{\alph{enumi})}
\tightlist
\item
  对每个 \(s \in G\)，\(U(s)\) 连续；b) 存在 e 的邻域 V，使得 \(U(V)\) 等度连续；c) 存在 E 中的总集 D，使得对每个 \(x \in D\)，映射 \(s \mapsto U(s)x\) 连续。
\end{enumerate}

这些条件显然是必要的。反之，假设条件 \(a\)、\(b\)、\(c\) 都满足。在 \(U(\mathrm{V})\) 上，逐点收敛拓扑与 D 中的逐点收敛拓扑相同（TVS, III, §3, No.~4, 命题 5）。因此映射 \((s,x)\mapsto U(s)x\) 从 \(\mathbf{V}\times \mathbf{E}\) 到 E 连续（GT, X, §2, No.~1, 命题 1 的推论 3）。由于 \(U(s_0s)x = U(s_0)(U(s)x)\) 对所有 \(s_0\in \mathbf{G}\)、\(s\in \mathbf{G}\)、\(x\in \mathbf{E}\) 都成立，可见 \(U\) 连续。

当 \(\mathbf{G}\) 局部紧时，条件 \(a)\) 和 \(b)\) 等价于以下条件：

\(a^{\prime})\) 对每个紧子集 \(\mathbf{K}\)（它属于 \(\mathbf{G}\)），\(U(\mathbf{K})\) 都等度连续。

\begin{enumerate}
\def\labelenumi{\arabic{enumi})}
\setcounter{enumi}{2}
\item
  假设 \(U\) 是 \(G\) 在 \(E\) 中的连续线性表示。对每个 \(s \in G\)，令 \(\widehat{U}(s)\) 为 \(U(s)\) 到 \(\widehat{E}\)（\(E\) 的完备化）的连续延拓。则 \(\widehat{U}\) 是 \(G\) 在 \(\widehat{E}\) 中的线性表示，满足注 2 的条件 \(a)\) 和 \(c)\)，并且根据 GT, X, §2, No.~2, 命题 4 也满足条件 \(b)\)。因此 \(\widehat{U}\) 是 \(G\) 在 \(\widehat{E}\) 中的连续线性表示。
\item
  当 \(\mathbf{E}\) 是赋范空间时，若 \(U\) 满足 \(\| U(s)\| = 1\) 对每个 \(s\in \mathbf{G}\) 成立，则称 U 为等距的。只要 \(\| U(s)\| \leqslant 1\) 对所有 \(s\in \mathbf{G}\) 成立就足够了，因为此时有
\end{enumerate}

\[
1 = \| 1 \| \leqslant \| U (s) \| \cdot \| U (s ^ {- 1}) \|,
\]

从而 \(\| U(s)\| = \| U(s^{-1})\| = 1\) 对所有 \(s\in G\) 成立。

命题 1. --- 若 G 是局部紧群且 E 是桶状空间，则 G 在 E 中的每个分别连续线性表示 U 都是连续的。

对 G 的每个紧子集 K，\(U(\mathrm{K})\) 关于逐点收敛拓扑是紧的（注 1），因而是等度连续的（TVS, III, §4, No.~2, Th. 1）；然后应用注 2。

引理 1. --- 设 G 为局部紧群，\(\rho\) 为 G 上取有限值且下半连续的数值函数，满足 \(\geqslant 0\)，并且满足 \(\rho(st) \leqslant \rho(s)\rho(t)\) 对所有 \(s, t \in \mathbf{G}\)。则 \(\rho\) 在 G 的每个紧子集上有上界。

存在 G 的非空开子集 U，使得 \(\rho\) 在 U 上有上界（GT, IX, §5, No.~4, Th. 2）。令 K 为 G 的紧子集。则 K 可被有限个集合 \(s_{1}U,\ldots,s_{n}U\) 覆盖。对每个 \(x\in U\)，有 \(\rho(s_{i}x)\leqslant\rho(s_{i})\rho(x)\)，因此 \(\rho\) 在各 \(s_{i}U\) 上有上界，从而在 K 上有上界。

引理 2. --- 设 G 为拓扑群，U 为 G 在赋范空间 E 中的线性表示，A 为 E 的稠密子集。假设对每个 \(s \in G\)，\(U(s)\) 连续，并且对每个 \(x \in A\)，从 G 到 E 的映射 \(s \mapsto U(s)x\) 连续。则 G 上的函数 \(s \mapsto g(s) = \|U(s)\|\) 下半连续，并满足 \(g(st) \leqslant g(s)g(t)\)。

令 B 为 E 的单位球。则 \(g(s) = \sup_{x \in \mathbf{B} \cap \mathbf{A}} \|U(s)x\|\)，且每个函数 \(s \mapsto \|U(s)x\|\) 在 G 上连续，因此 g 下半连续。另一方面，

\[
g (s t) = \| U (s) U (t) \| \leqslant \| U (s) \| \cdot \| U (t) \| = g (s) g (t).
\]

命题 2. --- 设 G 为局部紧群，U 为 G 在赋范空间 E 中的线性表示，A 为 E 的稠密子集。假设对每个 \(s \in G\)，\(U(s)\) 连续，并且对每个 \(x \in A\)，从 G 到 E 的映射 \(s \mapsto U(s)x\) 连续。则 U 连续。

因为根据引理 1 和 2，\(\|U(s)\|\) 在 G 的每个紧子集上有界，然后应用注 2 即可。

